\section{Planck constants and the determinantal structure of action}
\label{sec:action}
\subsection{Determinantal action scale and decimal valuation}

The absolute scale and the relative correction originate in different
objects. The former uses the local block of the dodecaphasic register
with step three,
\begin{equation}
 H_4=\begin{pmatrix}
 1&1&1&1\\1&1&-1&-1\\1&-1&1&-1\\1&-1&-1&1
 \end{pmatrix},\qquad
 G_H=\mathbb Z^4/H_4\mathbb Z^4.
 \label{eq:el-h4}
\end{equation}
The operation uses a local quotient. The three blocks of the complete
register are not replaced by a tensorization that would multiply
the number of sheets of this readout again.

\begin{lemma}[Number of sheets of the local quotient]
The Smith normal form of \(H_4\) is
\(\operatorname{diag}(1,2,2,4)\). Consequently,
\[
 G_H\simeq\mathbb Z/2\mathbb Z\oplus\mathbb Z/2\mathbb Z
 \oplus\mathbb Z/4\mathbb Z,\qquad |G_H|=16.
\]
\end{lemma}
\begin{proof}
The greatest common divisor of the entries is one. The minors of order two
are even, and one has value \(-2\), so their greatest common divisor
is two. Since \(H_4H_4^{\mathsf T}=4I\) and \(\det H_4=-16\), the
cofactors are \(\pm4\); the third determinantal divisor is four.
The fourth is \(16\). Successive ratios of the divisors
\(1,2,4,16\) are \(1,2,2,4\). Their product gives the quotient order.
\end{proof}

\begin{definition}[Determinantal action readout]
The multiplicative norm of the scalar section \(\alpha\), constant
on the sheets of \(G_H\), and one application of self-scaling
\(\varphi\) produce
\begin{equation}
 \Sigma_H=\varphi\,\mathrm N_{G_H}(\alpha),\qquad
 \mathrm N_{G_H}(\alpha)=\prod_{g\in G_H}\alpha=\alpha^{16}.
 \label{eq:el-lector-determinantal}
\end{equation}
The associated decimal chart determines
\(\operatorname{ord}_{10}(\Sigma_H)=\lfloor\log_{10}\Sigma_H\rfloor\).
\end{definition}

The sixteenth power thus arises from the number of sheets of the quotient,
once this local readout is fixed. The factor \(\varphi\) acts once on
the basis of the action line. The Smith calculation alone would not
select an arbitrary functional on the quotient: the multiplicative norm
and self-scaling form part of the declared operation.

\begin{theorem}[Decade of the action section]
The readout \eqref{eq:el-lector-determinantal}, applied to the HMT
coordinates already generated, satisfies
\begin{equation}
 10^{-34}<\Sigma_H<10^{-33},\qquad
 \operatorname{ord}_{10}(\Sigma_H)=-34.
 \label{eq:el-decada}
\end{equation}
\end{theorem}
\begin{proof}
The following rational intervals, much coarser than the available
cylinders, suffice:
\[
 \frac{1618}{1000}<\varphi<\frac{1619}{1000},\qquad
 \frac{7297}{10^6}<\alpha<\frac{7298}{10^6}.
\]
The function \((x,a)\mapsto xa^{16}\) is increasing in both positive
arguments. The two integer inequalities
\[
 1618\cdot7297^{16}>10^{65},\qquad
 1619\cdot7298^{16}<10^{66}
\]
imply
\[
 10^{-34}
 <\frac{1618}{1000}\left(\frac{7297}{10^6}\right)^{16}
 <\Sigma_H
 <\frac{1619}{1000}\left(\frac{7298}{10^6}\right)^{16}
 <10^{-33}.
\]
The definition of decimal order completes the proof. No SI unit
or reference mass enters these inequalities.
\end{proof}

The normalized section begins with
\(\Sigma_H=1.0462438381\ldots\times10^{-34}\). Its mantissa is not
identified with the action coordinate before or after return.
The determinantal readout fixes the decimal order; the two return
coordinates are obtained through the following readouts.


\subsection{Action sections and the dimensionless return quotient}

Let \(\mathcal L_S\) be an oriented action line and \(\mathcal U_S\)
its positive basis. The fifth-order character and the regional return
term have the expressions
\begin{align}
 H_5={}&\frac12\sqrt{\frac{1000\alpha}{\varphi}}-\alpha
   +\frac9{16}\alpha^2-\frac59\alpha^3
   +\frac7{48}\alpha^4-\frac1{54}\alpha^5,
   \label{eq:el-h5}\\
 D_A={}&\exp\!\left(-\frac{100\pi\alpha}{9}\right),
   \label{eq:el-da}\\
 \mathcal C_\pi(\alpha)={}&169\alpha^6+
       \frac{D_A\alpha^7}{1-D_A\alpha},
   \label{eq:el-cpi}\\
 \eta_{\mathrm{ret}}={}&H_5-\frac{90}{\pi}\mathcal C_\pi(\alpha).
   \label{eq:el-eta}
\end{align}
These are subsequent readouts of dodecaphasic closure and its
regional incidence. The integer \(169\) is retained as the regional
selector already constructed, and the coefficients of \(H_5\) as those of its
local character; they are not replaced by an interpolation of the Planck
constant. The formulas make explicit exactly which readouts are needed
for the electron stage. There is no need to introduce other
functionals of mass theory here.

For \(0<\alpha<1\), one has \(0<D_A<1\), so
\(D_A\alpha<1\). The rational term in \eqref{eq:el-cpi} preserves the
complete sum of the geometric tail:
\begin{equation}
 \frac{D_A\alpha^7}{1-D_A\alpha}
 =\sum_{n=0}^{\infty}D_A^{n+1}\alpha^{n+7}.
 \label{eq:el-cola-regional}
\end{equation}
It is not a truncation that may be silently modified when precision
changes. After the term of index \(N\), the remaining tail is
\(D_A^{N+2}\alpha^{N+8}/(1-D_A\alpha)\), allowing evaluation to
be controlled without recourse to a target quantity.

\begin{definition}[Action sections and return readout]
The sections before and after return are
\begin{equation}
 \hbar_{\mathrm{pre}}=H_5\,10^{-34}\mathcal U_S,
 \qquad
 \hbar_{\mathrm{ret}}=\eta_{\mathrm{ret}}\,10^{-34}\mathcal U_S.
 \label{eq:el-secciones-accion}
\end{equation}
Whenever \(\eta_{\mathrm{ret}}>0\), define
\begin{equation}
 \delta_{\mathrm{ret}}=H_5-\eta_{\mathrm{ret}},\qquad
 R_{\mathrm{act}}=\frac{\hbar_{\mathrm{pre}}}{\hbar_{\mathrm{ret}}}
 =\frac{H_5}{\eta_{\mathrm{ret}}}.
 \label{eq:el-ratio}
\end{equation}
\end{definition}

\begin{proposition}[Positivity and invariance of the ratio]
On the intervals used in the proof of \eqref{eq:el-decada},
\(0<\eta_{\mathrm{ret}}<H_5\), and therefore
\(R_{\mathrm{act}}>1\). The ratio does not change when the basis
\(\mathcal U_S\) is replaced by any other common positive basis.
\end{proposition}
\begin{proof}
The terms of \(\mathcal C_\pi\) are positive. For an elementary
bound, it suffices to use \(\alpha<0.008\), \(\varphi<1.619\),
\(\alpha>0.007\), \(D_A<1\), and \(\pi>3\). The root term in
\eqref{eq:el-h5} exceeds one, while the sum of the negative
terms is less than \(0.008001\); in particular, \(H_5>0.99\).
Likewise,
\[
 0<\frac{90}{\pi}\mathcal C_\pi
 <30\left(169(0.008)^6+
           \frac{(0.008)^7}{1-0.008}\right)<10^{-6}.
\]
Thus \(\eta_{\mathrm{ret}}>0.989999>0\) and
\(\eta_{\mathrm{ret}}<H_5\). Finally, the basis change divides
both action coordinates by the same positive factor, which cancels in
their quotient.
\end{proof}

Subsequent evaluation gives
\begin{align*}
 H_5&=1.0545718176461544696\ldots,\\
 \eta_{\mathrm{ret}}&=1.0545718169150390611\ldots,\\
 R_{\mathrm{act}}&=1.0000000006932817630\ldots.
\end{align*}
The difference and quotient are two readouts of the same sections,
one additive and one multiplicative. They do not represent two independent
Planck constants. The passage from angular action to action per complete
winding is expressed, in each section, by
\(h=2\pi\hbar\). The subsequent chart
\(\mathcal U_S\mapsto\mathrm{J\,s}\) assigns a unit, but does not
act retroactively on the decimal order already proved.

Regional return may also be written in an angular chart.
This reparametrization is not an indispensable additional dependency:
\eqref{eq:el-h5}--\eqref{eq:el-eta} directly supply
\(H_5\), \(\eta_{\mathrm{ret}}\), and \(R_{\mathrm{act}}\). If
angle and counterangle are introduced, both must be transported in the same
chart; degrees and radians are not mixed within a difference.

\subsubsection{Action character and provenance of its coefficients}
\label{sec:revision-planck}

The reduced Planck constant, \(\hbar\), is the physical object
motivating the realization of the angular action section. Its position in
the electronic dependency must be specified at two levels. The determinantal
norm \(\Sigma_H\) determines the decade of the scale. The character
\(H_5\) and regional return then determine the coordinates of the
action sections. This separation presents the result
required for the electron without anticipating the complete developments of
other sectors of dimensional mechanics.

The provenance of the fifth-order character can be made explicit through
invariants already fixed on the support. In the central chart we consider
\[
 B_c=\begin{pmatrix}7&2\\2&7\end{pmatrix},
 \qquad \lambda_+=9,\quad\lambda_-=5.
\]
The calendar has length \(12\), the step-\(3\) orbit has
length \(4\), and the relevant census quotient is
\(104976/1944=54\). The character coefficients satisfy
\begin{equation}
 \frac9{16}=\frac{\lambda_+}{4^2},\quad
 \frac59=\frac{\lambda_-}{\lambda_+},\quad
 \frac7{48}=\frac{\operatorname{tr}(B_c)/2}{4\cdot12},\quad
 \frac1{54}=\frac{1944}{104976}.
 \label{eq:revision-h5-coeficientes}
\end{equation}
The assignment of these invariants to degrees \(2,3,4,5\), with the
alternating signs of \eqref{eq:el-h5}, belongs to the definition of the
analytic character used. Making it explicit is essential: the set of
cardinalities, considered without the character and without the order of its degrees,
would admit other expressions. The constructive content lies in the
composition of the support, orientation and declared analytic rule.

\subsubsection{Regional continuation and renewal equation}

The selector \(169\) determines the degree-six impulse. Subsequent
prolongation is organized by the transfer \(D_A\) and a unit
normalization of the determinantal line. These rules admit a formulation
in formal series before substituting \(z=\alpha\), with \(D_A\) already
constructed. Write
\[
 \mathscr C(z)=169z^6+R(z),\qquad R(z)\in z^7\mathbb R[[z]].
\]
Concatenating a prolongation increases the degree by one and
multiplies its weight by \(D_A\). The first strict prolongation has
weight \(D_A\). Its renewal equation is therefore
\begin{equation}
 R(z)=D_Az^7+D_AzR(z).
 \label{eq:revision-renovacion}
\end{equation}

\begin{proposition}[Formal uniqueness of the normalized continuation]
Equation \eqref{eq:revision-renovacion} determines a unique series, and
\[
 \mathscr C(z)=169z^6+\frac{D_Az^7}{1-D_Az}
             =169z^6+\sum_{n=7}^{\infty}D_A^{n-6}z^n.
\]
The realization is analytic on \(|D_Az|<1\).
\end{proposition}
\begin{proof}
The element \(1-D_Az\) has constant term one and is invertible
in the formal power series ring. Solving the equation produces the displayed
series. Equivalently, the degree-seven coefficient is \(D_A\)
and each subsequent coefficient is \(D_A\) times its predecessor. The
geometric series converges absolutely in the indicated disk.
\end{proof}

Normalization of the first prolongation is an indispensable operational
datum. If only the degree-six impulse and
concatenation ratio were preserved, the series
\[
 169z^6+\lambda\frac{D_Az^7}{1-D_Az}
\]
would still share both data for any \(\lambda\).
The unit rule fixes \(\lambda=1\) before evaluation.
This locates each rule's contribution to uniqueness of the
readout, and the rational expression preserves the complete continuation.

\subsubsection{Reduced Planck constant and return of the action section}

In the dimensional chart of \eqref{eq:el-secciones-accion}, the
pre-return section is the model's construction intended for comparison
with the reduced Planck constant:
\[
 \hbar_{\mathrm{pre}}=H_5\,10^{-34}\mathcal U_S.
\]
The subsequent section also preserves the regional compensation
\(90\mathscr C(\alpha)/\pi\). Both belong to the same action
line, and their quotient \(R_{\mathrm{act}}\) measures the multiplicative
effect of return. Passage to \(h\) through \(2\pi\) uses the relation
between angular parametrization and a complete winding.

This construction effectively enters
\(R_{\mathrm{act}}^{80-54\alpha+6\Delta_4}\). The electron preserves
its central coordinate and trajectory; the chronological scale does not
replace that structure with an undifferentiated quantum. Nor is a mass
already expressed in MeV multiplied again by \(10^{-34}\):
the decade belongs to the action section, cancels in its dimensionless
quotient and reappears only where the corresponding dimensional
realization acts.

The numerical comparison of \(H_5\), the returned section and
\(h/(2\pi)\) is presented in Section~\ref{sec:revision-planck-contraste}.
This distinguishes the name of the physical object, the equation the
model assigns to it and the precision with which that equation reproduces it.

